\documentclass[12pt, a4paper, oneside]{report}
% ============================================================
%  Shared preamble — ViKi BS thesis (Innopolis University)
%  Compile with: xelatex (twice) — Cyrillic requires XeLaTeX
% ============================================================

\usepackage{fontspec}
\usepackage{polyglossia}
\setmainlanguage{english}
\setotherlanguage{russian}

\setmainfont{FreeSerif}[Ligatures=TeX]
\setsansfont{FreeSans}[Ligatures=TeX]
\setmonofont{DejaVu Sans Mono}[Scale=0.85]
\newfontfamily\cyrillicfont{FreeSerif}[Ligatures=TeX]
\newfontfamily\cyrillicfontsf{FreeSans}[Ligatures=TeX]
\newfontfamily\cyrillicfonttt{DejaVu Sans Mono}[Scale=0.85]

\usepackage{geometry}
\geometry{a4paper, top=25mm, bottom=25mm, left=30mm, right=15mm}

\usepackage{amsmath, amssymb, amsfonts, bm}
\usepackage{graphicx}
\usepackage{booktabs}
\usepackage{longtable}
\usepackage{array}
\usepackage{multirow}
\usepackage{caption}
\usepackage{subcaption}
\usepackage{enumitem}
\usepackage{xcolor}
\usepackage{setspace}
\usepackage{tikz}
\usetikzlibrary{arrows.meta, positioning, shapes.geometric, fit, backgrounds}
\usepackage[hidelinks]{hyperref}
\usepackage{tocloft}

\onehalfspacing
\setlength{\parindent}{1.25em}
\setlength{\parskip}{0.25em}

% --- IEEE-ish caption style -------------------------------
\captionsetup[figure]{labelfont=bf, labelsep=period, name=Fig., font=small}
\captionsetup[table]{labelfont=bf, labelsep=newline, justification=centering,
                     name=TABLE, font=small}
\renewcommand{\thetable}{\Roman{table}}

% --- placeholder macro ------------------------------------
% Every unresolved number / empirical claim is wrapped in \ph{}
% so it can be found with a single grep before submission.
\definecolor{phcol}{RGB}{176,0,32}
\newcommand{\ph}[1]{\textcolor{phcol}{\textbf{#1}}}
\newcommand{\TODO}[1]{\textcolor{phcol}{\small\textsf{[TODO: #1]}}}

% --- figure placeholder box -------------------------------
\newcommand{\figplaceholder}[2][0.62\textwidth]{%
  \begin{center}
  \fbox{\parbox[c][#2][c]{#1}{\centering\textcolor{phcol}{\textsf{\small
  [placeholder — insert figure]}}}}
  \end{center}}

% --- math notation ----------------------------------------
\newcommand{\R}{\mathbb{R}}
\newcommand{\SE}{\mathrm{SE}(3)}
\newcommand{\so}{\mathfrak{se}(3)}
\newcommand{\bq}{\bm{q}}
\newcommand{\bT}{\bm{T}}
\newcommand{\be}{\bm{e}}
\newcommand{\bx}{\bm{x}}
\newcommand{\bw}{\bm{w}}
\newcommand{\bJ}{\bm{J}}
\newcommand{\norm}[1]{\left\lVert #1 \right\rVert}
\DeclareMathOperator{\Log}{Log}
\DeclareMathOperator*{\argmin}{arg\,min}

% --- misc --------------------------------------------------
\newcommand{\viki}{ViKi}
\renewcommand{\cftsecleader}{\cftdotfill{\cftdotsep}}


\begin{document}

% ==========================================================
% TITLE PAGE
% ==========================================================
\begin{titlepage}
\centering
\vspace*{10mm}
{\large INNOPOLIS UNIVERSITY}\\[2mm]
{\normalsize Bachelor of Science in Computer Science}\\[30mm]

{\LARGE\bfseries ViKi: Vision-Based Kinematic Imitation\\[3mm]
for Object-Centric Retargeting of Human\\[3mm]
Manipulation Demonstrations to Robot Manipulators\par}
\vspace{15mm}

{\large Bachelor's Graduation Thesis}\\[25mm]

\begin{tabular}{rl}
\textit{Author:} & \ph{Artyom [Surname]}\\[4mm]
\textit{Supervisor:} & \ph{[Full Name] Kornaev}\\[4mm]
\textit{Group:} & \ph{[group]}\\
\end{tabular}

\vfill
{\normalsize Innopolis, Russia\\ \the\year}
\end{titlepage}

\pagenumbering{roman}

% ==========================================================
% ABSTRACT (EN)
% ==========================================================
\chapter*{Abstract}
\addcontentsline{toc}{chapter}{Abstract}

Imitation learning for robot manipulation is bottlenecked by the cost of
demonstration data: teleoperation requires dedicated hardware and an
operator per episode, which limits both throughput and the diversity of
recorded behaviour. Passive capture of bare-handed human demonstrations
removes that bottleneck, but existing passive pipelines either discard the
hand entirely and plan over sparse object keyframes, or reconstruct a full
parametric hand mesh from monocular video and retarget it in two decoupled
stages---smoothing the observed trajectory, solving inverse kinematics, then
smoothing the joint output again. The latter pattern is unsound: pre-filtering
in task space cannot guarantee smoothness of the nonlinear inverse-kinematics
solution, and it attenuates genuine fast motion together with perception noise.
Optimisation-based retargeting that fuses data fidelity, temporal smoothness
and joint limits into a single cost functional is by now standard---but only in
\emph{online teleoperation} systems, where a human remains in the loop and no
dataset is produced. This thesis closes that gap. It presents \viki, an offline
pipeline that captures bare-handed manipulation demonstrations with a
metrically calibrated multi-camera RGB-D rig anchored to a ChArUco workspace
frame, represents each demonstration as a dense object-relative wrist
trajectory in $\SE$ with a binary gripper state, and retargets it onto a robot
manipulator by minimising a single weighted cost functional that combines a
perception-confidence-weighted data term, a velocity-smoothness penalty and an
acceleration penalty subject to configuration, velocity and collision
constraints. Synthesised trajectories are replayed on the physical manipulator
before the dataset is finalised, so that kinematically infeasible episodes are
removed at the source rather than surfacing as policy failures later. The
system was evaluated on a constrained class of tabletop pick-and-place and
pouring tasks with a UR3 manipulator, two Azure Kinect DK cameras and one Intel
RealSense D435i. The proposed formulation reduced wrist-pose tracking error by
\ph{XX\%} and improved joint-space smoothness by \ph{XX\%} relative to a
sequential smooth--IK--smooth baseline, while replay screening rejected
\ph{XX\%} of episodes as infeasible; policies trained on the resulting
LeRobot-compatible dataset reached \ph{XX\%} success. The results indicate
that the online-teleoperation retargeting formulation transfers to the offline
dataset-generation regime, and that closing the embodiment gap by physical
replay before export is a practical alternative to filtering after training.

\vspace{3mm}
\noindent\textbf{Keywords:} imitation learning; motion retargeting; inverse
kinematics; RGB-D calibration; object-centric representation; robot learning
from human video.

% ==========================================================
% ABSTRACT (RU)
% ==========================================================
\chapter*{Аннотация}
\addcontentsline{toc}{chapter}{Аннотация}
\begin{russian}

Обучение роботов манипуляции методом имитации ограничено стоимостью получения
демонстраций: телеоперация требует специализированного оборудования и
оператора на каждый эпизод, что ограничивает и производительность записи, и
разнообразие поведения. Пассивный захват демонстраций, выполняемых голой рукой
человека, снимает это ограничение, однако существующие пассивные пайплайны либо
полностью отбрасывают кисть и планируют по разреженным ключевым кадрам объекта,
либо восстанавливают полную параметрическую mesh-модель кисти из монокулярного
видео и выполняют ретаргетинг в две развязанные стадии: сглаживание наблюдаемой
траектории, решение обратной задачи кинематики, повторное сглаживание в
пространстве суставов. Последняя схема методологически некорректна:
предварительная фильтрация в пространстве задачи не гарантирует гладкости
нелинейного решения обратной кинематики и подавляет истинное быстрое движение
вместе с шумом восприятия. Постановка ретаргетинга как единого функционала
стоимости, объединяющего точность соответствия данным, временную гладкость и
ограничения на суставы, к настоящему моменту стала стандартом --- но только в
системах \emph{онлайн-телеоперации}, где человек остаётся в контуре управления и
датасет не формируется. Настоящая работа закрывает этот разрыв. Предлагается
\viki\ --- офлайн-пайплайн, который захватывает демонстрации манипуляции голой
рукой при помощи метрически откалиброванного мультикамерного RGB-D рига,
привязанного к системе координат рабочей области по маркерной доске ChArUco,
представляет каждую демонстрацию как плотную объектно-относительную траекторию
запястья в $\SE$ с бинарным состоянием захвата и выполняет ретаргетинг на
роботизированный манипулятор путём минимизации единого взвешенного функционала,
объединяющего член соответствия данным со взвешиванием по достоверности
восприятия, штраф на скорость и штраф на ускорение при ограничениях на
конфигурацию, скорость и столкновения. Синтезированные траектории
воспроизводятся на физическом манипуляторе до финализации датасета, благодаря
чему кинематически нереализуемые эпизоды отсеиваются у источника, а не
проявляются позднее как отказы политики. Система оценена на ограниченном классе
настольных задач захвата--переноса и переливания с манипулятором UR3, двумя
камерами Azure Kinect DK и одной Intel RealSense D435i. Предложенная постановка
снизила ошибку отслеживания позы запястья на \ph{XX\%} и повысила гладкость в
пространстве суставов на \ph{XX\%} относительно последовательной базовой схемы
smooth--IK--smooth, при этом отбор по результатам воспроизведения отклонил
\ph{XX\%} эпизодов как нереализуемые; политики, обученные на полученном
LeRobot-совместимом датасете, достигли успешности \ph{XX\%}. Результаты
показывают, что формулировка ретаргетинга, разработанная для
онлайн-телеоперации, переносима в режим офлайн-генерации датасета, а закрытие
разрыва эмбодиментов через физическое воспроизведение до экспорта является
практичной альтернативой фильтрации после обучения.

\vspace{3mm}
\noindent\textbf{Ключевые слова:} обучение с имитацией; ретаргетинг движения;
обратная кинематика; калибровка RGB-D; объектно-центричное представление;
обучение роботов по видео человека.
\end{russian}

% ==========================================================
\chapter*{Acknowledgements}
\addcontentsline{toc}{chapter}{Acknowledgements}
\ph{[Placeholder.]} The author thanks \ph{[supervisor]} for supervision and
guidance throughout this work, and the Innopolis Robotics Society for
maintaining a parallel development fork of the system, whose independent
implementation attempts informed several of the architectural decisions
reported here.

\tableofcontents
\listoffigures
\listoftables

\cleardoublepage
\pagenumbering{arabic}
\setcounter{page}{1}

\chapter{Introduction}
\label{ch:intro}

\section{Background}
\label{sec:background}

Visuomotor policies such as Action Chunking Transformers (ACT) and Diffusion
Policy have made behaviour cloning a practical route to robot manipulation, but
their performance scales with the quantity and quality of demonstration data.
The dominant way of acquiring that data is teleoperation: a human operator
drives the robot through the task while joint states and camera streams are
recorded. Teleoperation yields action labels that are, by construction, exactly
executable on the target embodiment, which is its principal advantage. Its
disadvantages are equally structural. Each episode consumes an operator's time
in real time, the required hardware---a VR controller, an exoskeleton or a
haptic device---is expensive and embodiment-specific, and the resulting motion
is shaped as much by the latency and ergonomics of the interface as by the task
itself.

Human demonstrations recorded passively, with the demonstrator using their own
hands and no interface at all, invert this trade-off. Data collection becomes
cheap, fast and natural; a demonstrator can record dozens of episodes in the
time a teleoperator records one, and the recorded motion reflects how the task
is actually performed rather than how it can be performed through a control
interface. What is lost is precisely the property teleoperation guarantees:
there are no action labels, and the kinematics of a human arm and hand do not
correspond to the kinematics of a robot manipulator. Bridging that gap is the
retargeting problem, and it is the subject of this thesis.

\section{Problem Statement}
\label{sec:problem}

Passive capture makes demonstration data cheap but produces observations, not
actions. Converting an observed human hand trajectory into a robot joint
trajectory that (i) reproduces the task-relevant motion, (ii) is smooth enough
for a policy to imitate, and (iii) is physically executable on the target
manipulator, is an underdetermined problem contaminated by perception noise.

Two failure modes dominate. The first is noise amplification through the
kinematic chain. Monocular or neural depth estimates of hand keypoints are
noisy at the centimetre scale; inverse kinematics is nonlinear, so a bounded
task-space perturbation may map to an unbounded joint-space excursion near
singularities. The widespread response---low-pass filtering the observed
trajectory before solving inverse kinematics, then filtering the joint
trajectory afterwards---does not solve this. Task-space pre-filtering is
applied in the wrong space to guarantee joint-space smoothness, and it removes
genuine high-frequency task content along with the noise. This is the
garbage-in-garbage-out pattern that motivates the present work.

The second failure mode is representational. A trajectory recorded as an
absolute sequence of positions in a camera or world frame is tied to the exact
spatial layout of the recording session. If the object is moved, the trajectory
is wrong; if the embodiment changes, it is wrong in a different way. What
generalises is the relation between the hand and the object being manipulated,
not the absolute path of the hand through space.

\section{Research Gap}
\label{sec:gap}

Both failure modes have known remedies, and the literature has adopted them---
but in disjoint communities.

Optimisation-based retargeting, in which data fidelity, temporal smoothness and
joint limits are posed as a single constrained optimisation over the robot
configuration, is standard in \emph{online teleoperation}. AnyTeleop and the
associated \texttt{dex-retargeting} library established the canonical form;
Open-TeleVision, Bunny-VisionPro and ACE build directly on it; DexFlow extends
the smoothness penalty to second order. In every case, however, a human remains
in the control loop, the optimisation runs at interactive rates, and no dataset
is produced.

Conversely, \emph{offline} pipelines that generate datasets from passive human
video adopt object-centric representations but not the joint-optimisation
formulation. ORION plans over sparse object keyframes and does not perform
per-frame hand retargeting at all. OKAMI reconstructs a parametric body-and-hand
model and retargets in two decoupled stages, warping a reference trajectory to
a new object position and then solving inverse kinematics---structurally the
same sequential pattern criticised in Section~\ref{sec:problem}.

A targeted prior-art review conducted for this work (Chapter~\ref{ch:lit},
Section~\ref{sec:positioning}) found no published system combining passive
markerless capture, a metrically calibrated multi-camera RGB-D rig, an
object-centric dense trajectory representation, single-cost-functional
retargeting, imitation-learning-ready dataset export, and physical replay used
as a pre-export validation step. The nearest systems each satisfy some criteria
and miss others. The gap addressed here is therefore not the invention of a new
retargeting operator, but the transfer of an established formulation into a
regime where it has not been applied, together with the sensing and validation
infrastructure that regime requires.

\section{Research Questions and Hypotheses}
\label{sec:rq}

The work is guided by one primary research question and three sub-questions.

\vspace{2mm}
\noindent\textbf{RQ.} \emph{How can human manipulation demonstrations captured
passively by a calibrated multi-camera RGB-D rig be retargeted onto a robot
manipulator so as to yield datasets suitable for training visuomotor imitation
policies?}

\begin{enumerate}[label=\textbf{RQ\arabic*.}, leftmargin=2.2em]
  \item How does formulating retargeting as a single weighted cost functional,
        combining a confidence-weighted data term with velocity and
        acceleration penalties, affect the fidelity and smoothness of the
        resulting joint trajectories compared with a sequential
        smooth--IK--smooth pipeline?
  \item What accuracy can be achieved for wrist pose estimation by
        confidence-weighted fusion of multiple calibrated RGB-D views anchored
        to a common ChArUco workspace frame, and how does the residual error
        decompose across its independent sources?
  \item To what extent does replaying synthesised trajectories on the physical
        manipulator, prior to dataset export, improve the executability of the
        exported dataset?
\end{enumerate}

\noindent The corresponding hypotheses are stated below. Following the
convention for experimental computer-science hypotheses, each identifies an
independent variable, a dependent variable, a population and a relation.

\begin{description}[leftmargin=3.2em, style=nextline]
  \item[\textnormal{\textbf{H1 (behaviour).}}]
  Retargeting posed as a single weighted cost functional produces joint
  trajectories with lower task-space tracking error and higher smoothness than
  a sequential smooth--IK--smooth pipeline, on the same recorded demonstrations.
  \emph{Independent variable:} retargeting formulation (joint vs.\ sequential).
  \emph{Dependent variables:} wrist-pose tracking RMSE; spectral arc length of
  the joint trajectory. \emph{Population:} recorded demonstrations of the
  constrained task class defined in Section~\ref{sec:tasks}.
  \emph{Relation:} the joint formulation dominates on both measures.

  \item[\textnormal{\textbf{H2 (behaviour).}}]
  Weighting the data term by per-frame perception confidence reduces tracking
  error relative to uniform weighting, and the reduction is larger on frames
  where hand landmarks are partially occluded.
  \emph{Independent variable:} presence of confidence weighting.
  \emph{Dependent variable:} wrist-pose tracking RMSE, stratified by occlusion.
  \emph{Population:} as above. \emph{Relation:} monotone reduction, amplified
  under occlusion.

  \item[\textnormal{\textbf{H3 (dependability).}}]
  Screening trajectories by physical replay before export increases the
  proportion of exported episodes that execute successfully on the manipulator.
  \emph{Independent variable:} presence of the replay screening stage.
  \emph{Dependent variable:} executable fraction of the exported dataset.
  \emph{Population:} as above. \emph{Relation:} positive.
\end{description}

The questions were checked against the FINER criteria. They are
\emph{feasible}, in that the hardware, solver and evaluation protocol are all
available and the task class is deliberately constrained; \emph{interesting}
and \emph{relevant} to the robot-learning community, because demonstration cost
is a recognised bottleneck; \emph{novel} in the specific sense established in
Section~\ref{sec:gap}, namely a regime transfer rather than a new operator; and
\emph{ethical}, as the study involves no human subjects beyond the author
performing demonstrations and no personally identifying data is retained.

\section{Contributions}
\label{sec:contributions}

This thesis makes the following contributions.

\begin{enumerate}[leftmargin=2em]
  \item A retargeting formulation for the \emph{offline} setting that poses
        data fidelity, velocity smoothness and acceleration smoothness as one
        constrained optimisation over the robot configuration, with the data
        term weighted by per-frame perception confidence and made robust by a
        Huber loss (Section~\ref{sec:retarget}).
  \item An object-centric dense trajectory representation that stores each
        demonstration both relative to the manipulated object and relative to a
        fixed ChArUco workspace anchor, permitting downstream consumers to
        select the representation that generalises better
        (Section~\ref{sec:objcentric}).
  \item A multi-camera RGB-D capture and calibration procedure with
        confidence-weighted depth fusion and an explicit treatment of
        cross-device timestamp alignment, together with an error budget that
        separates the three independent contributions to residual pose error
        (Sections~\ref{sec:calib}--\ref{sec:fusion}, Section~\ref{sec:errbudget}).
  \item A physical-replay validation stage that filters kinematically
        infeasible trajectories before dataset export, and an empirical
        characterisation of what that stage rejects
        (Sections~\ref{sec:replay}, \ref{sec:res-replay}).
  \item An open implementation exporting LeRobot-compatible datasets, evaluated
        end-to-end by training ACT and Diffusion Policy on the generated data
        (Chapter~\ref{ch:impl}).
\end{enumerate}

\section{Thesis Structure}
\label{sec:structure}

Chapter~\ref{ch:lit} reviews prior work on hand representation, retargeting
formulations, task representation, multi-camera RGB-D calibration and dataset
generation, and positions \viki\ against the nearest published systems.
Chapter~\ref{ch:method} describes the design of the system and the derivation
of the retargeting cost functional. Chapter~\ref{ch:impl} reports the
implementation, the evaluation protocol and the experimental results.
Chapter~\ref{ch:disc} interprets those results, decomposes the error budget and
states the limitations. Chapter~\ref{ch:concl} concludes and outlines future
work.

\chapter{Literature Review}
\label{ch:lit}

\section{Preamble}
\label{sec:lit-preamble}

This chapter surveys the state of knowledge relevant to converting passively
recorded human manipulation demonstrations into robot-executable datasets. It
is organised not chronologically but by the architectural elements of the
proposed system, so that for every design decision taken in
Chapter~\ref{ch:method} the alternatives considered and the reason for the
choice are visible.

Section~\ref{sec:lit-hand} examines how the human hand is represented, from
full parametric meshes to the minimal wrist-pose-plus-gripper abstraction.
Section~\ref{sec:lit-retarget} contrasts sequential and single-cost-functional
retargeting formulations and reviews the solvers used to realise them.
Section~\ref{sec:lit-repr} compares object-centric and world-frame task
representations. Section~\ref{sec:lit-calib} covers multi-camera RGB-D
calibration, depth fusion and cross-device synchronisation, with emphasis on
reported quantitative accuracy. Section~\ref{sec:lit-dataset} reviews dataset
formats, downstream policies and demonstration curation.
Section~\ref{sec:positioning} positions \viki\ against the nearest published
systems, and Section~\ref{sec:lit-conclusion} states the resulting
methodological choices.

Sources were drawn from arXiv (\texttt{cs.RO}, \texttt{cs.CV}, \texttt{cs.LG}),
the proceedings of CoRL, RSS, ICRA, IROS, CVPR and NeurIPS, peer-reviewed
sensor and calibration literature, and vendor documentation for the depth
cameras and solvers employed. Where a source is a recent preprint whose
peer-review status could not be confirmed, this is stated explicitly.

\section{Representing the Human Hand}
\label{sec:lit-hand}

\subsection{Full parametric reconstruction}

A large body of work recovers the complete kinematics of the demonstrator.
OKAMI \cite{okami} combines SMPL-H for the body with MANO for the hand and
retargets the two separately. DexCanvas \cite{dexcanvas} builds a human--robot
correspondence dataset by fitting MANO to data from twenty-two infrared motion
capture cameras and two RGB-D sensors. DexMV \cite{dexmv} likewise recovers a
full MANO hand together with object pose from video before translating the
demonstration into simulation.

Full reconstruction yields a rich representation---the pose of every finger
joint---and is necessary when the target embodiment is a multi-fingered
dexterous hand. It also introduces a heavy perception stack and an additional,
poorly characterised error source: the mesh fit itself. For a parallel-jaw
gripper, most of that representational richness is discarded downstream.

\subsection{Minimal representation: wrist pose and gripper state}

The alternative, adopted in this work, represents the hand as a wrist pose in
$\SE$ together with a scalar or binary gripper state. This is not a concession
made for tractability; it is the representation used by the mainstream of
current teleoperation systems. Open-TeleVision \cite{opentv} maps the operator's
wrist pose to the arm end-effector through closed-loop inverse kinematics and
encodes a parallel gripper by the single thumb--index vector. Bunny-VisionPro
\cite{bunny} and ACE \cite{ace} adopt the same decomposition, ACE additionally
demonstrating that it transfers across dexterous hands, parallel grippers and
quadrupeds.

Methodological support for down-weighting the wrist comes from the ablation
study of Xin et al.\ \cite{keyobj}, who assemble a comprehensive retargeting
objective---global hand pose, overall hand shape, relative fingertip positions
and fingertip orientations---and remove terms one at a time. They find that
precise tracking of the global wrist pose is frequently unnecessary and can
degrade fingertip accuracy, and consequently replace the wrist-position
objective with a thumb-tip objective while regularising wrist orientation only
weakly. The implication for a wrist-plus-gripper representation is that the
wrist target should enter the optimisation as a weighted soft task rather than
a hard constraint.

\subsection{Egocentric and monocular capture}

Several systems reduce capture-hardware requirements further. EgoMimic
\cite{egomimic} records bare-handed demonstrations through Project Aria glasses
and co-trains on human and robot data. DexWild \cite{dexwild} collects
in-the-wild bare-hand data through a low-cost portable device and likewise
co-trains. A recent egocentric RGB-D system \cite{shadowing} generates
trajectories offline using analytic damped-least-squares inverse kinematics and
replays them on a low-cost arm. These works confirm a trend toward lighter
capture rigs, but each accepts monocular or single-view geometry, and none
combines it with the metric multi-view accuracy that the present work requires.

\section{Retargeting Formulations}
\label{sec:lit-retarget}

\subsection{The single cost functional}

The prevailing modern formulation poses retargeting as a constrained nonlinear
least-squares or quadratic program over the robot configuration $\bq_t$ at each
timestep, minimising a weighted sum of a data-fidelity term and a temporal
smoothness term subject to joint limits. AnyTeleop \cite{anyteleop} gives the
canonical instance: the objective sums, over keypoints, the squared residual
between a scaled human keypoint vector and the corresponding robot
forward-kinematics vector, plus a first-order penalty on configuration change
between consecutive frames. The accompanying \texttt{dex-retargeting} library
exposes a vector optimiser for real-time teleoperation, a position optimiser
explicitly recommended for offline dataset processing, and a contact-aware
DexPilot optimiser; it is built on Pinocchio and solved with sequential
quadratic programming.

Two refinements are directly relevant here. First, per-term confidence
weighting: in the formulation reported by Open-TeleVision \cite{opentv}, each
keypoint residual is scaled by a distance-dependent switching weight that
increases when a pinch is detected. The mechanism is precisely analogous to
weighting by a per-frame perception confidence, and it establishes that
weighting individual data terms by a scalar signal is accepted practice.
Second, second-order smoothness: DexFlow \cite{dexflow} augments the standard
objective with a dynamic smoothing weight and Hessian regularisation to enforce
$C^2$ continuity, penalising curvature and not merely velocity. Robustification
by a Huber loss on the data term, reported in recent dexterous teleoperation
work \cite{dexteleop0}, provides a mechanism for rejecting perception outliers
inside the optimiser rather than smoothing them away beforehand.

\subsection{The sequential pattern and its critique}

Against this, the sequential pattern---filter the observed task-space
trajectory, solve inverse kinematics per frame, filter the joint output---
remains common in offline video-to-robot pipelines. Its deficiency is
structural rather than empirical. Filtering is applied in task space, whereas
the smoothness requirement is on the joint trajectory; because inverse
kinematics is nonlinear, a smooth Cartesian input does not imply a smooth joint
output, particularly near singularities. Pre-filtering also attenuates genuine
fast motion indiscriminately with noise. In the single-functional formulation
the smoothness penalty acts directly on $\bq_t$, so joint-space smoothness is
obtained by construction.

\subsection{Offline pipelines and their retargeting operators}

ORION \cite{orion} does not perform per-frame hand retargeting: it constructs an
object-centric plan over sparse keyframes from an open-world object graph.
OKAMI \cite{okami} warps a reference trajectory to a new object position and
then solves inverse kinematics---the sequential pattern. HOWTransfer
\cite{howtransfer}, the closest published pipeline in workflow terms, performs
passive markerless capture and transfers trajectories using flow-matching
$\SE$ grasp sampling with Laplacian trajectory editing, not a single weighted
cost functional; its representation is hand-centric rather than object-centric.
No offline pipeline located in this review adopts the AnyTeleop-style single
functional.

\subsection{Solvers}

Pinocchio underlies the retargeting stack of AnyTeleop, Open-TeleVision and
ACE. Pink \cite{pink} exposes weighted inverse kinematics directly as a
quadratic program: each task contributes a quadratic residual with a weight
matrix normalising task coordinates to common units, while joint limits,
velocity limits and self-collision barriers enter as inequality constraints,
optionally with control-barrier formulations. This structure admits the entire
proposed cost functional natively---a frame task for the wrist target, posture
and velocity regularisation tasks for smoothness, and limits as constraints---
solved as one program per frame. The solver choice made in
Chapter~\ref{ch:method} therefore lies on the mainstream path rather than
diverging from it.

\section{Task Representation}
\label{sec:lit-repr}

Object-centric representations dominate recent work on cross-environment and
cross-embodiment generalisation. ORION \cite{orion} represents a task as the
three-dimensional object trajectory relative to initial and final object
positions, arguing that this generalises across visual backgrounds, camera
shifts, spatial layouts and novel object instances, and that full-body
reconstruction of the demonstrator is unnecessary. SPOT \cite{spot} encodes each
task as the $\SE$ object-pose trajectory relative to the target explicitly in
order to decouple embodiment-specific actions from sensory input, enabling
learning from action-free human demonstrations, and conditions a diffusion
policy on the object trajectory. Flow-based formulations \cite{novaflow} make an
equivalent argument while relaxing the rigid-object assumption that
pose-trajectory methods require.

OKAMI \cite{okami} occupies an intermediate position: the representation is
demonstrator-centric, but object-aware warping adapts trajectories to new object
positions---an implicit acknowledgement that pure world-frame replay fails when
the object moves. Vi-RoIL \cite{viroil} is object-centric by construction,
expressing wrist trajectories and keyframe hand poses in the object frame, but
recovers metric scale from a monocular view of a marker cube and evaluates only
in simulation.

The trade-off is that object-centric representations require robust object pose
or segmentation, and pose-trajectory variants assume rigidity, whereas
world-frame trajectories are simpler but tied to the recording layout. Storing
both, as proposed in Section~\ref{sec:objcentric}, defers the choice to the
downstream consumer at negligible storage cost.

\section{Multi-Camera RGB-D Calibration, Fusion and Synchronisation}
\label{sec:lit-calib}

\subsection{Fiducial choice}

ChArUco boards combine the sub-pixel corner accuracy of a checkerboard with the
robust identification and occlusion tolerance of ArUco markers. A $6\times8$
board provides up to thirty-five corners per view against four for a single
AprilTag, and remains detectable under partial occlusion and at the image
boundary where a plain checkerboard fails. Reported reprojection errors for
ChArUco and ArUco calibration are typically below half a pixel. AprilTags are
easier to detect at long range and oblique incidence and suffice for
extrinsic-only estimation, but are ill-suited to intrinsic calibration owing to
the small number of corners. Current practice is therefore to use ChArUco for
intrinsics and either fiducial for extrinsics, preferring AprilTags for oblique
overhead views.

\subsection{Reported accuracy of multi-Azure-Kinect rigs}

Romeo et al.\ \cite{romeo} provide the cleanest direct comparison of calibration
procedures for multi-Kinect rigs, reporting that three-dimensional procedures
outperform two-dimensional chessboard approaches, with mean distances between
corresponding point-cloud points of $21.4$~mm at colour resolution and
$9.9$~mm at infrared resolution for a static scene, and $20.9$~mm and $7.4$~mm
respectively for a dynamic scene; body-joint alignment error was $35.4$~mm.
Darwish et al.\ \cite{darwish} calibrate five Azure Kinects against a
ChArUco-coded cube, combining reprojection error with point-to-plane and
patch-to-plane terms to compensate depth noise. Kim et al.\ \cite{kimcal}
perform coarse registration from body tracking followed by feature-match
refinement, retaining above ninety-five per cent inliers. Bu and Lee
\cite{bulee} use markerless coarse calibration with coloured ICP refinement and
report---importantly for rig geometry---that ICP fails for near-opposing camera
placements owing to insufficient point-cloud overlap.

Single-sensor specifications bound the achievable accuracy from below. The
Azure Kinect DK is specified at a random error standard deviation of at most
$17$~mm and a typical systematic error below $11$~mm plus $0.1$ per cent of
distance \cite{mskinect}; independent evaluation confirms these figures
\cite{tolgyessy}, and comparative work finds improved spatial and temporal
accuracy relative to Kinect~v2 beyond $2.5$~m \cite{kurillo}. The Intel
RealSense D435i is specified at depth accuracy below one per cent of distance
\cite{intel}. The practical consequence is that centimetre-level, not
millimetre-level, agreement should be expected between calibrated views.

\subsection{Depth fusion}

Naive averaging of overlapping depth observations is inferior to
confidence-weighted fusion. Weighting by distance and by the angle between the
surface normal and the viewing direction---so that surfaces observed near
perpendicular contribute more than obliquely observed ones---is established
practice \cite{fusionpatent}. Learned confidence maps estimated from the depth
signal itself are used to re-weight depth features before fusion
\cite{zeng}, and surfel-based reconstruction attaches an explicit confidence
value to each element and performs curvature-weighted registration
\cite{hrbf}. Both sensor families used in this work expose per-pixel
invalidation or confidence signals that can drive such weights directly.

\subsection{Cross-device synchronisation}

Azure Kinect devices support hardware synchronisation through daisy-chained or
star-wired trigger cabling, with subordinate devices required to start before
the master and depth exposures staggered to prevent laser interference; the
vendor documentation states explicitly that equal timestamps between depth and
colour streams, or between cameras, are not guaranteed \cite{mskinect}.
RealSense devices provide an analogous inter-camera synchronisation mode, with
the counter-intuitive property that timestamps appear aligned when hardware
synchronisation is disabled and exhibit bounded drift when it is enabled, the
drift itself serving as verification \cite{intel, ridgerun}.

The pitfall for a heterogeneous rig is that hardware clocks on different device
families are not mutually comparable, since each counts its own cycles from its
own epoch. Documented practice for such rigs is to stamp every arriving frame
with a single host-side monotonic clock and to align streams by nearest
timestamp in that common domain, accepting millisecond-scale precision
\cite{ridgerun}. Because Azure Kinect and RealSense triggers are mutually
incompatible, this is the only available mechanism for aligning the two
families, and it is the mechanism adopted in Section~\ref{sec:sync}.

\subsection{Hand--eye calibration}

The OpenCV \texttt{calibrateHandEye} routine implements the Tsai--Lenz,
Park--Martin, Horaud--Dornaika, Andreff and Daniilidis solvers. A comparative
noise study \cite{alihandeye} finds the Daniilidis, Chou and Park methods most
robust to translation noise, with Daniilidis yielding the best translation
estimates, while Tsai, Li and Lu degrade as rotation noise grows. Methods that
decouple rotation from translation are faster but propagate rotation error into
the translation estimate; Daniilidis solves both simultaneously through dual
quaternions. Because the present system anchors cameras and robot to a shared
ChArUco workspace frame, classical hand--eye estimation is largely bypassed;
where an explicit camera-to-robot transform is required, Daniilidis is used.

\section{Dataset Generation and Downstream Policies}
\label{sec:lit-dataset}

The LeRobot HDF5 conventions together with ACT and Diffusion Policy constitute
the default imitation-learning toolchain and define the export target of this
work. OpenVLA-OFT \cite{openvlaoft} reports that parallel decoding, action
chunking, a continuous action representation and $L1$ regression raise average
success across four task suites from $76.5$ to $97.1$ per cent while increasing
action-generation throughput twenty-six-fold, and is therefore a credible
alternative consumer of the generated data.

Curation matters as much as generation. Demo-SCORE \cite{demoscore} trains a
classifier on policy rollouts to filter low-quality demonstrations and reports
an absolute success-rate improvement exceeding fifteen to thirty-five per cent
over training on all demonstrations. Smoothness-based scoring of trajectories
has been reported to yield comparable gains at a fraction of the data
\ph{[citation to be confirmed]}. This motivates treating dataset curation as an
explicit pipeline stage rather than an afterthought, and---in the present
design---moving it upstream of export in the form of physical replay screening.

\section{Positioning of the Proposed System}
\label{sec:positioning}

Table~\ref{tab:positioning} compares \viki\ against the nearest published
systems along the six criteria that jointly define the proposed architecture:
(C1) passive markerless capture without teleoperation hardware; (C2)
multi-camera metrically calibrated RGB-D sensing; (C3) object-centric dense
trajectory representation; (C4) retargeting as a single confidence-weighted
cost functional; (C5) export of an imitation-learning-ready dataset; and (C6)
physical replay used to validate the dataset before export.

\begin{table}[htbp]
\caption{Comparison of the Proposed System Against the Nearest Published Work\\
Along the Six Criteria Defined in Section~\ref{sec:positioning}}
\label{tab:positioning}
\centering
\small
\begin{tabular}{lcccccc}
\toprule
System & C1 & C2 & C3 & C4 & C5 & C6 \\
\midrule
\viki\ (proposed)      & \checkmark & \checkmark & \checkmark & \checkmark & \checkmark & \checkmark \\
HOWTransfer \cite{howtransfer} & \checkmark & $\sim$ & --- & --- & $\sim$ & \checkmark \\
Vi-RoIL \cite{viroil}          & $\sim$ & --- & \checkmark & --- & $\sim$ & --- \\
ORION \cite{orion}             & \checkmark & $\sim$ & \checkmark & --- & --- & --- \\
OKAMI \cite{okami}             & \checkmark & $\sim$ & $\sim$ & --- & $\sim$ & --- \\
EgoMimic \cite{egomimic}       & \checkmark & --- & --- & --- & $\sim$ & --- \\
DexWild \cite{dexwild}         & $\sim$ & --- & --- & --- & $\sim$ & --- \\
Tool-as-Interface \cite{toolif}& \checkmark & $\sim$ & $\sim$ & --- & $\sim$ & --- \\
RoboPaint \cite{robopaint}     & --- & \checkmark & $\sim$ & --- & \checkmark & \checkmark \\
DexCanvas \cite{dexcanvas}     & --- & \checkmark & --- & --- & $\sim$ & --- \\
\bottomrule
\end{tabular}

\vspace{1mm}
{\footnotesize \checkmark~= satisfied; $\sim$~= partially satisfied;
---~= not satisfied.}
\end{table}

No located system satisfies all six criteria simultaneously. HOWTransfer
\cite{howtransfer} is closest in workflow, satisfying C1 and C6, but is
hand-centric rather than object-centric, uses calibrated stereo RGB rather than
metric RGB-D, and does not employ a single cost functional. Vi-RoIL
\cite{viroil} is closest in representation, satisfying C3, but relies on a
monocular view of a marker cube and is evaluated only in simulation. The
intersection of the online joint-optimisation retargeting literature with the
offline object-centric dataset-generation literature is therefore unoccupied,
and it is that intersection which this work targets.

Several of the systems listed are recent preprints
\cite{howtransfer, robopaint, shadowing, dexcanvas} whose peer-review status
could not be confirmed at the time of writing; their reported results are
treated as provisional.

\section{Conclusion of the Review}
\label{sec:lit-conclusion}

The review supports the following methodological choices, which
Chapter~\ref{ch:method} implements.

\begin{enumerate}[leftmargin=2em]
  \item \emph{Hand representation.} A wrist pose in $\SE$ with a binary gripper
        state is adopted, following \cite{opentv, bunny, ace}, and the wrist
        target enters as a weighted soft task following the ablation evidence
        of \cite{keyobj}. Parametric mesh reconstruction is rejected as
        unnecessary for a parallel-jaw target and as an additional error source.
  \item \emph{Retargeting.} A single weighted cost functional is adopted,
        following \cite{anyteleop, dexflow}, with confidence weighting after
        \cite{opentv} and Huber robustification after \cite{dexteleop0}. The
        sequential smooth--IK--smooth pattern is rejected on the structural
        grounds set out in Section~\ref{sec:lit-retarget}.
  \item \emph{Solver.} Pink on Pinocchio \cite{pink} is adopted, since it
        expresses the entire functional natively as a constrained quadratic
        program and is the solver family used by the reference teleoperation
        systems.
  \item \emph{Task representation.} Object-relative trajectories are adopted
        following \cite{orion, spot}, and stored alongside workspace-anchored
        trajectories so that the choice is deferred downstream.
  \item \emph{Sensing.} ChArUco workspace anchoring is adopted; centimetre-level
        cross-camera agreement is assumed on the evidence of \cite{romeo};
        camera placement preserves substantial view overlap on the evidence of
        \cite{bulee}; depth is fused with distance, incidence-angle and sensor
        confidence weights following \cite{fusionpatent, zeng}.
  \item \emph{Synchronisation.} The two Azure Kinects are hardware
        synchronised; the RealSense, which cannot share their trigger, is
        aligned through a single host-side monotonic clock following
        \cite{ridgerun}.
  \item \emph{Curation.} Dataset curation is treated as an explicit stage and
        moved upstream of export in the form of physical replay screening,
        motivated by \cite{demoscore} and by the replay-based validation of
        \cite{howtransfer}.
\end{enumerate}

\chapter{Design and Methodology}
\label{ch:method}

\section{System Overview}
\label{sec:overview}

\viki\ is an offline pipeline: demonstrations are recorded, processed and
exported in distinct stages, with no requirement that any stage run at
interactive rates. This is a deliberate design choice. Relaxing the real-time
constraint that shapes teleoperation systems permits global optimisation over
whole trajectories, iterative refinement, and a validation stage that executes
on physical hardware---none of which is available to an online system.

The pipeline comprises six stages, shown in Fig.~\ref{fig:pipeline}.

\begin{enumerate}[leftmargin=2em]
  \item \emph{Capture.} Synchronised RGB-D streams are recorded from a
        calibrated multi-camera rig while a demonstrator performs the task
        bare-handed.
  \item \emph{Perception.} Hand landmarks are detected per view, lifted to
        three dimensions using sensor depth, and fused across views with
        confidence weights to yield a wrist pose in $\SE$ and a gripper state.
  \item \emph{Representation.} The wrist trajectory is expressed both relative
        to the manipulated object and relative to the ChArUco workspace anchor.
  \item \emph{Retargeting.} A robot joint trajectory is obtained by minimising
        a single weighted cost functional subject to kinematic and collision
        constraints.
  \item \emph{Replay validation.} The synthesised trajectory is executed on the
        physical manipulator; proprioception is recorded and infeasible
        episodes are rejected.
  \item \emph{Export.} Surviving episodes are written in a LeRobot-compatible
        format for downstream policy training.
\end{enumerate}

\begin{figure}[htbp]
\centering
\begin{tikzpicture}[
  node distance=4mm,
  box/.style={draw, rounded corners=1pt, align=center, font=\footnotesize,
              minimum height=8mm, text width=22mm, inner sep=2pt},
  arr/.style={-{Latex[length=2mm]}, thick}]
\node[box] (cap) {Capture\\\scriptsize 2$\times$Kinect + RealSense};
\node[box, right=of cap] (per) {Perception\\\scriptsize landmarks + fusion};
\node[box, right=of per] (rep) {Representation\\\scriptsize object-relative $\SE$};
\node[box, below=8mm of rep] (ret) {Retargeting\\\scriptsize single cost functional};
\node[box, left=of ret] (val) {Replay\\\scriptsize physical validation};
\node[box, left=of val] (exp) {Export\\\scriptsize LeRobot / HDF5};
\draw[arr] (cap) -- (per);
\draw[arr] (per) -- (rep);
\draw[arr] (rep) -- (ret);
\draw[arr] (ret) -- (val);
\draw[arr] (val) -- (exp);
\draw[arr] (val.north) -- ++(0,4mm) -| node[pos=0.25, above, font=\scriptsize]
  {reject / re-solve} (ret.east);
\end{tikzpicture}
\caption{Stages of the \viki\ pipeline. Capture, perception and representation
are embodiment-agnostic; retargeting, replay validation and export are
specific to the target manipulator. The feedback edge indicates that episodes
failing replay are either re-solved with modified weights or discarded.}
\label{fig:pipeline}
\end{figure}

\section{Hardware and Workspace}
\label{sec:hardware}

The capture rig comprises two Azure Kinect DK cameras and one Intel RealSense
D435i observing a tabletop workspace, with a UR3 manipulator mounted at a fixed
position relative to that workspace. The Azure Kinect was selected for the
primary views because its time-of-flight depth sensing provides substantially
better range estimation than neural regression from RGB, and because its
specified error characteristics---random error standard deviation at most
$17$~mm and systematic error below $11$~mm plus $0.1$ per cent of distance
\cite{mskinect}---are documented and independently verified \cite{tolgyessy}.
The RealSense provides a supplementary close-range view at higher depth
accuracy \cite{intel}.

The two Kinects are placed to preserve substantial overlap in their fields of
view. This is a direct consequence of the finding of Bu and Lee \cite{bulee}
that iterative-closest-point refinement fails between near-opposing views for
lack of point-cloud overlap; near-orthogonal placement retains both
triangulation baseline and overlap.

\section{Calibration}
\label{sec:calib}

\subsection{Intrinsics and extrinsics}

Intrinsic parameters for every colour and depth stream are estimated from
ChArUco board observations. Extrinsic calibration follows the workspace-anchor
principle: rather than expressing cameras relative to one another, each camera
is calibrated against a ChArUco board rigidly fixed to the workspace, defining
a common world frame $W$. For camera $k$ this yields
\begin{equation}
  \bT^{W}_{C_k} \in \SE,
  \label{eq:extrinsic}
\end{equation}
the transform from the camera frame $C_k$ to the workspace frame $W$. Anchoring
to the workspace rather than to a reference camera has two consequences: adding
or moving a camera does not invalidate the calibration of the others, and the
robot base can be registered into the same frame without a separate
camera-to-camera chain.

\subsection{Robot registration}

The manipulator base frame $R$ is registered to $W$ by observing a ChArUco
target held in a known end-effector pose across a set of configurations, giving
$\bT^{W}_{R}$ directly. Where an explicit camera-to-robot transform is required
instead, the Daniilidis dual-quaternion solver is used, following the
robustness evidence of \cite{alihandeye}, with fifteen to twenty poses spanning
a hemisphere.

\section{Synchronisation}
\label{sec:sync}

The two Azure Kinects are hardware-synchronised in master/subordinate
configuration, with subordinate devices started before the master and depth
exposures staggered to avoid mutual laser interference \cite{mskinect}. The
RealSense cannot share that trigger; its synchronisation mode is incompatible
at the electrical level.

Cross-family alignment therefore uses a single host-side monotonic clock. Every
frame, from every device, is stamped on arrival with
\begin{equation}
  \tau = \texttt{time.monotonic\_ns()},
  \label{eq:hostclock}
\end{equation}
and streams are aligned by nearest neighbour in $\tau$. Two points deserve
emphasis. First, per-device hardware timestamps are \emph{not} used for
cross-device arithmetic: distinct device families count their own cycles from
their own epochs, and the vendor documentation states explicitly that equal
timestamps across cameras are not guaranteed \cite{mskinect}. Second, a
monotonic clock rather than wall-clock time is required, since wall-clock time
may be adjusted discontinuously by time synchronisation daemons, which would
corrupt interval computations mid-recording.

Residual alignment jitter is measured rather than assumed
(Section~\ref{sec:res-calib}), and hardware synchronisation of the Kinect pair
is verified by confirming bounded drift in the inter-camera timestamp
difference \cite{ridgerun}.

\section{Perception and Confidence-Weighted Fusion}
\label{sec:fusion}

\subsection{Landmark detection and lifting}

Hand landmarks are detected in each colour stream, yielding for view $k$ at
time $t$ a set of image-plane landmarks with associated visibility scores
$v^{k}_{i,t} \in [0,1]$. Each landmark is lifted to three dimensions using the
registered depth map rather than a learned depth estimate, exploiting the
time-of-flight measurement directly, and transformed into the workspace frame
by \eqref{eq:extrinsic}.

\subsection{Fusion}

Estimates from multiple views are combined by confidence-weighted averaging
rather than by naive mean. For landmark $i$ at time $t$, the fused position is
\begin{equation}
  \hat{\bx}_{i,t} \;=\;
  \frac{\sum_{k} w^{k}_{i,t}\, \bx^{k}_{i,t}}
       {\sum_{k} w^{k}_{i,t}},
  \label{eq:fusion}
\end{equation}
with weights
\begin{equation}
  w^{k}_{i,t} \;=\;
  \underbrace{v^{k}_{i,t}}_{\text{detector}}
  \cdot
  \underbrace{c^{k}_{i,t}}_{\text{sensor}}
  \cdot
  \underbrace{\frac{1}{\left(d^{k}_{i,t}\right)^{2}}}_{\text{range}}
  \cdot
  \underbrace{\max\!\left(0, \cos\theta^{k}_{i,t}\right)}_{\text{incidence}},
  \label{eq:weights}
\end{equation}
where

\begin{center}
\begin{tabular}{@{\hspace{10mm}}p{16mm}@{}p{95mm}}
$v^{k}_{i,t}$ & detector visibility score for landmark $i$ in view $k$;\\
$c^{k}_{i,t}$ & sensor depth confidence or validity flag at that pixel;\\
$d^{k}_{i,t}$ & range from camera $k$ to the landmark;\\
$\theta^{k}_{i,t}$ & angle between the local surface normal and the viewing direction.
\end{tabular}
\end{center}

\noindent The range and incidence factors follow established practice in
geometric depth fusion, where surfaces observed near perpendicular are weighted
above obliquely observed ones because they are measured more accurately
\cite{fusionpatent}; the detector and sensor factors follow confidence-map
re-weighting \cite{zeng}. The product form is chosen so that a zero in any
factor---an undetected landmark, an invalid depth pixel, a grazing view---
removes that view's contribution entirely rather than merely attenuating it.

\subsection{Wrist pose and gripper state}

The wrist pose $\bT^{W}_{H,t} \in \SE$ is obtained by fitting an orthonormal
frame to the fused wrist and knuckle landmarks. The gripper state is derived
from the thumb--index distance,
\begin{equation}
  g_t \;=\;
  \mathbb{1}\!\left[\,\norm{\hat{\bx}_{\text{thumb},t}
  - \hat{\bx}_{\text{index},t}} < \tau_g \,\right],
  \label{eq:gripper}
\end{equation}
with hysteresis on $\tau_g$ to suppress chatter at the threshold. This single
scalar is the entire hand representation beyond the wrist frame, following the
convention established for parallel-jaw targets in \cite{opentv, ace}.

\section{Object-Centric Representation}
\label{sec:objcentric}

Let $\bT^{W}_{O,t}$ denote the pose of the manipulated object in the workspace
frame. Each demonstration is stored in two forms. The workspace-anchored form
is $\bT^{W}_{H,t}$ directly; the object-relative form is
\begin{equation}
  \bT^{O}_{H,t} \;=\; \left(\bT^{W}_{O,t}\right)^{-1} \bT^{W}_{H,t}.
  \label{eq:objrel}
\end{equation}

The object-relative form is what transfers. If the same task is to be performed
with the object at a new pose $\bT^{R}_{O'}$ expressed in the robot frame, the
desired end-effector pose at time $t$ is
\begin{equation}
  \bT^{R,\mathrm{des}}_{E,t} \;=\;
  \bT^{R}_{O'} \; \bT^{O}_{H,t} \; \bT^{H}_{E},
  \label{eq:transfer}
\end{equation}
where $\bT^{H}_{E}$ is a fixed offset relating the human wrist frame to the
robot end-effector frame, determined once per embodiment. Trajectories are
retained densely rather than reduced to keyframes, since the continuous
velocity profile of the demonstration carries task-relevant information---
approach speed, contact timing---that sparse keyframes discard.

Both representations are exported. Storing both costs little and defers to the
downstream consumer a choice that the literature does not settle universally:
object-relative trajectories generalise across object placement
\cite{orion, spot} but depend on object pose estimation, whereas
workspace-anchored trajectories are simpler and adequate within a fixed layout.

\section{Retargeting as a Single Cost Functional}
\label{sec:retarget}

\subsection{Formulation}

Let $\bq_t \in \R^{n}$ be the manipulator configuration at time $t$ and
$f(\bq_t) \in \SE$ the forward kinematics of the end-effector. The task-space
residual with respect to the desired pose \eqref{eq:transfer} is the
logarithmic map of the pose error,
\begin{equation}
  \be_t(\bq_t) \;=\;
  \Log\!\left(
    \left(\bT^{R,\mathrm{des}}_{E,t}\right)^{-1} f(\bq_t)
  \right)^{\vee} \in \R^{6}.
  \label{eq:residual}
\end{equation}

Retargeting is then posed as the single constrained problem
\begin{equation}
\begin{aligned}
  \min_{\bq_{1:T}} \;\; J(\bq_{1:T}) \;=\; \sum_{t=1}^{T} \Big[\;
   & \underbrace{\omega_t \,\rho_\delta\!\left(
       \norm{\be_t(\bq_t)}_{\bm{\Sigma}}\right)}_{\text{data fidelity}}
   \;+\;
     \underbrace{\lambda_v \norm{\bq_t - \bq_{t-1}}^{2}}_{\text{velocity}}
   \\[1mm]
   &+\;
     \underbrace{\lambda_a \norm{\bq_t - 2\bq_{t-1} + \bq_{t-2}}^{2}}
       _{\text{acceleration}}
   \;+\;
     \underbrace{\lambda_r \norm{\bq_t - \bq_{\mathrm{ref}}}^{2}}
       _{\text{posture}}
  \;\Big]
\end{aligned}
\label{eq:cost}
\end{equation}
subject to
\begin{align}
  \bq_{\min} \;\le\; &\bq_t \;\le\; \bq_{\max}, \label{eq:cjoint}\\
  \norm{\bq_t - \bq_{t-1}} \;&\le\; \dot{\bq}_{\max}\,\Delta t,
    \label{eq:cvel}\\
  h_j(\bq_t) \;&\ge\; 0, \quad j = 1,\dots,m, \label{eq:ccoll}
\end{align}
where

\begin{center}
\begin{tabular}{@{\hspace{10mm}}p{28mm}@{}p{85mm}}
$\omega_t$ & per-frame perception confidence weight, defined in \eqref{eq:conf};\\
$\rho_\delta$ & Huber loss with transition parameter $\delta$;\\
$\bm{\Sigma}$ & diagonal matrix normalising translational and rotational residual units;\\
$\lambda_v, \lambda_a, \lambda_r$ & weights on velocity, acceleration and posture regularisation;\\
$h_j$ & collision and self-collision barrier functions.
\end{tabular}
\end{center}

\subsection{Confidence weighting}

The per-frame weight aggregates the landmark confidences that produced the
wrist estimate at time $t$,
\begin{equation}
  \omega_t \;=\;
  \left(\frac{1}{|\mathcal{I}|}\sum_{i \in \mathcal{I}}
        \max_k w^{k}_{i,t}\right)^{\alpha},
  \label{eq:conf}
\end{equation}
with $\mathcal{I}$ the landmark subset defining the wrist frame and $\alpha$
controlling how sharply low-confidence frames are discounted. The construction
is the offline analogue of the distance-dependent switching weights used in
online teleoperation retargeting \cite{opentv}: frames in which the hand is
partially occluded or poorly observed exert proportionally less influence on
the solution, and the smoothness terms carry the trajectory through them
instead of the data term dragging it toward a spurious observation.

\subsection{Robustification}

The Huber loss $\rho_\delta$ bounds the influence of gross outliers---a
mis-detected hand, a depth dropout that survives the confidence weighting---at
the level of the data term itself, rather than allowing a filtering stage to
smear the outlier across neighbouring frames \cite{dexteleop0}. This is the
mechanism that replaces pre-smoothing: outliers are down-weighted where they
enter, not attenuated after they have already propagated.

\subsection{Why a single functional}

Equation~\eqref{eq:cost} differs from the sequential alternative in one
structurally important respect. The smoothness penalties act on $\bq_t$, in the
space where smoothness is required, and they are traded against data fidelity
by the solver rather than applied blindly beforehand. In the sequential
pipeline, task-space filtering cannot guarantee joint-space smoothness because
the forward kinematics is nonlinear, and it cannot distinguish task-relevant
fast motion from noise because it operates without reference to the data term.
Here the trade-off is explicit: on high-confidence frames $\omega_t$ is large
and the solution tracks the observation; on low-confidence frames the
smoothness terms dominate and the solution interpolates.

Second-order smoothness, expressed by the acceleration term, follows the
$C^2$-continuity argument of DexFlow \cite{dexflow}. A trajectory that is
first-order smooth may still exhibit curvature discontinuities that appear as
jerk on the physical manipulator and as label noise to a policy trained on the
data.

\subsection{Solution}

The problem is solved with Pink on Pinocchio \cite{pink}. Each term of
\eqref{eq:cost} maps onto a weighted task in the quadratic program---a frame
task for the data term, posture and velocity regularisation tasks for the
smoothness terms---while \eqref{eq:cjoint}--\eqref{eq:ccoll} enter as
inequality constraints, with collision avoidance expressed through control
barrier functions. Levenberg--Marquardt damping stabilises the solution near
singularities.

Because the setting is offline, the problem is solved over the whole trajectory
rather than frame by frame, so the acceleration term is well defined at every
interior timestep and the solution is not constrained by the causality
requirement that an online solver faces. The trajectory is initialised from the
previous solution warm-started along $t$, and boundary conditions at $t=1,2$
are handled by replicating the initial configuration.

\subsection{Interpolation}

Where gaps remain after fusion---landmarks lost across all views---the
trajectory is completed by cubic spline interpolation over $\SE$ prior to
optimisation. Dynamic filters that impose a motion model, such as Kalman or
extended Kalman filters, are deliberately not applied: they introduce a prior
about the dynamics that is subsequently re-estimated by the optimiser, and the
interaction between the two priors is difficult to characterise. The optimiser
already imposes the required regularity through
$\lambda_v$ and $\lambda_a$.

\section{Replay Validation}
\label{sec:replay}

Synthesised trajectories are executed on the physical manipulator before the
dataset is finalised. The stage serves three purposes.

First, it filters infeasibility. A trajectory may satisfy
\eqref{eq:cjoint}--\eqref{eq:ccoll} in the kinematic model and still fail on
hardware through controller tracking limits, unmodelled collisions or workspace
boundary conditions. Such episodes are rejected rather than exported.

Second, it closes the embodiment gap at the source. The proprioceptive states
recorded during replay are those the robot actually attained, not those the
optimiser predicted. Exporting the recorded rather than the synthesised states
means the dataset contains honest observation--action pairs from the target
embodiment, which is precisely the property that teleoperated datasets possess
by construction and that naive video-derived datasets lack.

Third, it provides a curation signal. Episodes are ranked by tracking residual
during replay, and the ranking is available to downstream consumers as a
quality score, in the spirit of rollout-based demonstration curation
\cite{demoscore} but applied before rather than after policy training.

Failure of an episode triggers either rejection or re-solution of
\eqref{eq:cost} with modified weights---typically an increased $\lambda_v$ or a
relaxed data weight in the region of failure---as indicated by the feedback
edge in Fig.~\ref{fig:pipeline}. The number of re-solution attempts is bounded
to prevent overfitting the trajectory to the replay controller.

\section{Dataset Export}
\label{sec:export}

Surviving episodes are written in a LeRobot-compatible HDF5 layout. Each
episode contains the synchronised camera streams, the recorded joint states and
gripper states from replay, the object-relative and workspace-anchored wrist
trajectories, per-frame confidence weights, and the replay tracking residual.
Retaining the confidence weights and residuals in the exported artefact allows
downstream curation to be repeated with different thresholds without
re-recording.

\section{Assumptions and Simplifications}
\label{sec:assumptions}

The following assumptions delimit the scope of the work and are revisited in
Section~\ref{sec:limitations}.

\begin{enumerate}[leftmargin=2em]
  \item \emph{Parallel-jaw target.} The hand is reduced to a wrist pose and a
        binary gripper state. Multi-fingered dexterous targets would require
        the finger-level representation rejected in
        Section~\ref{sec:lit-hand}.
  \item \emph{Rigid objects.} The object-relative representation
        \eqref{eq:objrel} presupposes a rigid body with a well-defined pose.
        Deformable objects would require a flow-based representation
        \cite{novaflow}.
  \item \emph{Single-arm, quasi-static tasks.} The task class is restricted to
        single-arm tabletop manipulation at moderate speed, where the
        quasi-static assumption underlying purely kinematic retargeting holds.
  \item \emph{Fixed workspace.} The ChArUco anchor is rigidly attached to the
        workspace and the camera rig is not moved within a recording session.
  \item \emph{Kinematic-only transfer.} Forces and contact dynamics are not
        estimated. Tasks whose success depends on force regulation are outside
        the scope.
\end{enumerate}

\chapter{Implementation and Results}
\label{ch:impl}

\section{Software Architecture}
\label{sec:arch}

The system is implemented in Python without a robotics middleware layer.
Avoiding ROS was a deliberate choice: the pipeline is a batch data-processing
system rather than a distributed real-time control system, and the node graph,
message serialisation and build tooling that middleware provides would add
operational complexity without addressing any requirement of the offline
setting. Communication with the manipulator uses \texttt{ur-rtde} directly.

The implementation is organised into the layers listed in
Table~\ref{tab:modules}. The separation follows the pipeline stages of
Fig.~\ref{fig:pipeline}, with the coordinate-frame contract---every spatial
quantity carries an explicit frame identifier---enforced at module boundaries so
that frame errors surface as type errors rather than as silent numerical
mistakes.

\begin{table}[htbp]
\caption{Principal Modules of the Implementation}
\label{tab:modules}
\centering
\small
\begin{tabular}{p{32mm}p{45mm}p{45mm}}
\toprule
Module & Responsibility & Principal dependencies \\
\midrule
\texttt{capture} & Device drivers, hardware sync, host-clock stamping &
  Azure Kinect SDK, \texttt{pyrealsense2} \\
\texttt{calibration} & ChArUco detection, intrinsics, extrinsics, robot
  registration & OpenCV \\
\texttt{perception} & Landmark detection, depth lifting, confidence-weighted
  fusion & MediaPipe, OpenCV \\
\texttt{representation} & Object pose, object-relative transform, spline
  interpolation & NumPy, SciPy \\
\texttt{retargeting} & Cost functional assembly and solution & Pink, Pinocchio \\
\texttt{replay} & Trajectory execution, proprioception logging, screening &
  \texttt{ur-rtde} \\
\texttt{export} & LeRobot/HDF5 serialisation & h5py, LeRobot \\
\texttt{interface} & Recording control and inspection & FastAPI, Streamlit \\
\bottomrule
\end{tabular}
\end{table}

The Azure Kinect depth engine requires a software rendering context inside the
containerised deployment; this is obtained by forcing software rasterisation
through the appropriate driver environment variables. The system is released
under the MIT licence.

\section{Experimental Setup}
\label{sec:setup}

\subsection{Hardware configuration}

Experiments used a UR3 manipulator with a parallel-jaw gripper, two Azure
Kinect DK cameras in master/subordinate hardware synchronisation and one Intel
RealSense D435i, all observing a tabletop workspace anchored by a ChArUco
board. Camera placement was near-orthogonal to preserve view overlap
(Section~\ref{sec:hardware}). Recording ran at \ph{XX}~Hz.

\subsection{Task suite}
\label{sec:tasks}

Evaluation used a deliberately constrained task class, chosen so that the
quasi-static and rigid-object assumptions of Section~\ref{sec:assumptions} hold
and so that success is unambiguously measurable:

\begin{enumerate}[leftmargin=2em]
  \item \emph{Pick-and-place.} Grasp a rigid object and place it at a marked
        target location.
  \item \emph{Pick-and-place with displacement.} As above, but with the object
        initialised at a pose not seen during demonstration, to test the
        object-relative representation.
  \item \emph{Pouring.} Grasp a container and tip it over a target, testing
        orientation fidelity rather than position alone.
  \item \ph{[Additional task, if included.]}
\end{enumerate}

For each task, \ph{N} demonstrations were recorded by \ph{one} demonstrator.

\subsection{Baselines and ablations}

The proposed formulation \eqref{eq:cost} is compared against:

\begin{enumerate}[leftmargin=2em]
  \item \textbf{B1 --- Sequential.} Task-space low-pass filtering, per-frame
        inverse kinematics, joint-space low-pass filtering. This is the
        pattern used by offline pipelines such as \cite{okami} and the
        principal comparison for H1.
  \item \textbf{B2 --- Joint, no confidence.} Equation~\eqref{eq:cost} with
        $\omega_t \equiv 1$. Isolates the contribution of confidence weighting
        for H2.
  \item \textbf{B3 --- Joint, no acceleration.} Equation~\eqref{eq:cost} with
        $\lambda_a = 0$. Isolates the contribution of second-order smoothness.
  \item \textbf{B4 --- Joint, no Huber.} Squared loss in place of
        $\rho_\delta$. Isolates outlier robustification.
  \item \textbf{B5 --- World-frame.} Workspace-anchored rather than
        object-relative targets, evaluated on the displaced-object task only.
\end{enumerate}

\section{Evaluation Protocol and Metrics}
\label{sec:metrics}

Metrics were fixed before the final experimental runs, so that the criteria of
success were not selected after observing the outcomes.

\subsection{Retargeting fidelity}

Task-space tracking error between the desired end-effector pose
\eqref{eq:transfer} and the achieved pose $f(\bq_t)$, reported separately for
translation and rotation:
\begin{equation}
  \mathrm{RMSE}_{\mathrm{pos}} =
  \sqrt{\frac{1}{T}\sum_{t} \norm{\bm{p}_t^{\mathrm{des}} - \bm{p}_t}^{2}},
  \qquad
  \mathrm{RMSE}_{\mathrm{rot}} =
  \sqrt{\frac{1}{T}\sum_{t} \norm{\Log(\bm{R}_t^{\mathrm{des}\,\top}
  \bm{R}_t)}^{2}}.
  \label{eq:rmse}
\end{equation}

\subsection{Smoothness}

Spectral arc length of the joint velocity profile, a dimensionless smoothness
measure that is invariant to amplitude and duration, together with mean
absolute joint jerk. Lower magnitudes indicate smoother motion.

\subsection{Feasibility}

The proportion of synthesised episodes that execute on hardware without
controller fault, joint limit violation or collision, and the mean tracking
residual reported by the controller during replay.

\subsection{Downstream policy success}

ACT and Diffusion Policy trained on the exported dataset and evaluated over
\ph{N} rollouts per task, reporting success rate under a task-specific success
predicate.

\subsection{Calibration and sensing accuracy}

Reprojection error of the ChArUco calibration in pixels; cross-camera
registration error as the mean distance between corresponding points in
overlapping regions; and residual timestamp misalignment between device
families.

\section{Calibration and Sensing Results}
\label{sec:res-calib}

\begin{table}[htbp]
\caption{Calibration and Synchronisation Accuracy}
\label{tab:calib}
\centering
\small
\begin{tabular}{lcc}
\toprule
Quantity & Measured & Reference from literature \\
\midrule
ChArUco reprojection error (px)        & \ph{X.XX} & $<0.5$ typical \\
Cross-camera registration, IR (mm)     & \ph{XX.X} & $7.4$--$9.9$ \cite{romeo} \\
Cross-camera registration, colour (mm) & \ph{XX.X} & $20.9$--$21.4$ \cite{romeo} \\
Kinect pair timestamp drift (ms/min)   & \ph{X.XX} & $<1$ \cite{ridgerun} \\
Kinect--RealSense residual offset (ms) & \ph{X.XX} & --- \\
Wrist position repeatability (mm)      & \ph{X.XX} & --- \\
\bottomrule
\end{tabular}
\end{table}

\TODO{Populate Table~\ref{tab:calib} from the calibration runs; discuss whether
measured cross-camera registration falls within the range reported by
\cite{romeo}, and if not, whether the discrepancy is attributable to board
size, view overlap or working distance.}

\section{Retargeting Fidelity}
\label{sec:res-fidelity}

\begin{table}[htbp]
\caption{Retargeting Fidelity and Smoothness, Proposed Formulation Versus
Sequential Baseline}
\label{tab:fidelity}
\centering
\small
\begin{tabular}{lcccc}
\toprule
Method & $\mathrm{RMSE}_{\mathrm{pos}}$ (mm) & $\mathrm{RMSE}_{\mathrm{rot}}$
(deg) & SAL & Jerk (rad/s$^3$) \\
\midrule
B1 Sequential            & \ph{XX.X} & \ph{X.X} & \ph{$-$X.XX} & \ph{XX.X} \\
Proposed \eqref{eq:cost} & \ph{XX.X} & \ph{X.X} & \ph{$-$X.XX} & \ph{XX.X} \\
\midrule
Relative change          & \ph{$-$XX\%} & \ph{$-$XX\%} & \ph{$-$XX\%}
                         & \ph{$-$XX\%} \\
\bottomrule
\end{tabular}
\end{table}

\figplaceholder{45mm}
\begin{figure}[htbp]
\caption{Joint trajectory for a representative pick-and-place episode under the
sequential baseline (B1) and the proposed formulation. \TODO{insert plot}}
\label{fig:traj}
\end{figure}

\TODO{Report whether H1 is supported. State the statistical test used and the
number of episodes; if the improvement is present on one metric but not the
other, say so plainly rather than reporting only the favourable one.}

\section{Ablation Studies}
\label{sec:res-ablation}

\begin{table}[htbp]
\caption{Ablation of the Individual Terms of the Cost Functional}
\label{tab:ablation}
\centering
\small
\begin{tabular}{llcc}
\toprule
Variant & Removed term & $\mathrm{RMSE}_{\mathrm{pos}}$ (mm) & SAL \\
\midrule
Proposed & ---                        & \ph{XX.X} & \ph{$-$X.XX} \\
B2 & confidence weighting $\omega_t$   & \ph{XX.X} & \ph{$-$X.XX} \\
B3 & acceleration penalty $\lambda_a$  & \ph{XX.X} & \ph{$-$X.XX} \\
B4 & Huber robustification $\rho_\delta$ & \ph{XX.X} & \ph{$-$X.XX} \\
\bottomrule
\end{tabular}
\end{table}

\begin{table}[htbp]
\caption{Effect of Confidence Weighting Stratified by Landmark Occlusion}
\label{tab:occlusion}
\centering
\small
\begin{tabular}{lcc}
\toprule
Frame subset & B2 (uniform) & Proposed (weighted) \\
\midrule
All frames                & \ph{XX.X} & \ph{XX.X} \\
Low-occlusion frames      & \ph{XX.X} & \ph{XX.X} \\
High-occlusion frames     & \ph{XX.X} & \ph{XX.X} \\
\bottomrule
\end{tabular}

\vspace{1mm}
{\footnotesize Values are $\mathrm{RMSE}_{\mathrm{pos}}$ in mm. Occlusion
stratification uses the median of $\omega_t$ as the threshold.}
\end{table}

\TODO{Report whether H2 is supported, and specifically whether the benefit of
confidence weighting is larger on high-occlusion frames as hypothesised. If the
effect is uniform across strata, the mechanism proposed in
Section~\ref{sec:retarget} is not the operative one and this must be stated.}

\section{Replay Feasibility}
\label{sec:res-replay}

\begin{table}[htbp]
\caption{Outcome of Physical Replay Screening}
\label{tab:replay}
\centering
\small
\begin{tabular}{lccc}
\toprule
Task & Episodes synthesised & Rejected at replay & Rejection cause (dominant) \\
\midrule
Pick-and-place              & \ph{NN} & \ph{NN} (\ph{XX\%}) & \ph{[cause]} \\
Pick-and-place (displaced)  & \ph{NN} & \ph{NN} (\ph{XX\%}) & \ph{[cause]} \\
Pouring                     & \ph{NN} & \ph{NN} (\ph{XX\%}) & \ph{[cause]} \\
\midrule
Total                       & \ph{NN} & \ph{NN} (\ph{XX\%}) & --- \\
\bottomrule
\end{tabular}
\end{table}

\TODO{Report the distribution of rejection causes---joint limit, singularity
proximity, collision, controller tracking fault---since the distribution
indicates which constraint in \eqref{eq:cjoint}--\eqref{eq:ccoll} is
insufficiently modelled rather than merely how many episodes failed.}

\section{Downstream Policy Performance}
\label{sec:res-policy}

\begin{table}[htbp]
\caption{Success Rate of Policies Trained on the Exported Dataset}
\label{tab:policy}
\centering
\small
\begin{tabular}{lcc}
\toprule
Task & ACT & Diffusion Policy \\
\midrule
Pick-and-place              & \ph{XX\%} & \ph{XX\%} \\
Pick-and-place (displaced)  & \ph{XX\%} & \ph{XX\%} \\
Pouring                     & \ph{XX\%} & \ph{XX\%} \\
\midrule
Mean                        & \ph{XX\%} & \ph{XX\%} \\
\bottomrule
\end{tabular}
\end{table}

\begin{table}[htbp]
\caption{Effect of Representation Choice on the Displaced-Object Task}
\label{tab:repr}
\centering
\small
\begin{tabular}{lc}
\toprule
Representation & Success rate \\
\midrule
B5 Workspace-anchored (world frame) & \ph{XX\%} \\
Object-relative \eqref{eq:objrel}   & \ph{XX\%} \\
\bottomrule
\end{tabular}
\end{table}

\TODO{Report whether H3 is supported by comparing policy success on datasets
with and without replay screening. Note that this requires training on the
unscreened dataset as well, which doubles the training budget; if that
comparison is not run, H3 must be evaluated on the executability metric alone
and the weaker claim stated explicitly.}

\chapter{Analysis and Discussion}
\label{ch:disc}

\section{Interpretation of the Results}
\label{sec:interpretation}

\TODO{This section interprets rather than restates Chapter~\ref{ch:impl}. Write
it after the numbers are in, structured as: what the key findings mean for each
hypothesis; which results support them and which do not; how the findings
compare with prior work; explanation of discrepancies; unexpected findings.}

The three hypotheses of Section~\ref{sec:rq} address separable mechanisms, and
the argument for each is independent of the others.

\paragraph{H1 --- formulation.} The claim is that posing retargeting as one
weighted functional dominates the sequential pipeline on both fidelity and
smoothness simultaneously. The distinction matters because the sequential
pipeline can be made arbitrarily smooth by increasing filter strength, at the
cost of fidelity; the interesting question is therefore not whether either
metric can be improved in isolation but whether the trade-off frontier itself
moves. \ph{[Insert result.]} If the proposed method improves one metric while
degrading the other, the correct conclusion is that the frontier has not moved
and the improvement is a re-parameterisation of the same trade-off.

\paragraph{H2 --- confidence weighting.} The mechanism proposed in
Section~\ref{sec:retarget} predicts a specific signature: the benefit of
weighting should concentrate on frames where perception is degraded, because
that is where uniform weighting drags the solution toward a bad observation.
A uniform improvement across occlusion strata would indicate that the weighting
is acting as a global regulariser rather than through the proposed mechanism.
\ph{[Insert result.]}

\paragraph{H3 --- replay screening.} The claim is about dependability rather
than accuracy: screening should raise the executable fraction of the exported
dataset. \ph{[Insert result.]} The stronger claim---that screening improves
downstream policy success---requires the paired training run noted in
Section~\ref{sec:res-policy} and should not be asserted on the executability
metric alone.

\section{Comparison with Prior Work}
\label{sec:comparison}

Direct numerical comparison with the systems reviewed in
Chapter~\ref{ch:lit} is not meaningful, since none is evaluated on the same
task suite, embodiment or sensing configuration. What can be compared is
qualitative behaviour and the class of guarantees provided.

Relative to OKAMI \cite{okami}, the present system discards the parametric
body-and-hand reconstruction and the two-stage warp-then-solve structure. The
expected consequence is a lower perception error floor---time-of-flight depth
in place of a mesh fit---at the cost of being restricted to parallel-jaw
targets. Relative to ORION \cite{orion}, the present system retains the dense
wrist trajectory that ORION discards, which should preserve velocity-profile
information at the cost of a dependency on reliable hand observation. Relative
to the online teleoperation systems \cite{anyteleop, opentv, ace}, the
formulation is the same but the operating regime is not: the offline setting
permits whole-trajectory optimisation with a well-defined acceleration term at
every interior timestep and admits a physical validation stage, neither of
which an interactive-rate solver can perform.

The measured cross-camera registration error
(Table~\ref{tab:calib}) should be compared against the $7.4$--$21.4$~mm range
reported by Romeo et al.\ \cite{romeo}. \TODO{State whether the measured value
falls in that range and, if not, diagnose the cause rather than reporting the
discrepancy without explanation.}

\section{Error Budget}
\label{sec:errbudget}

A central methodological commitment of this work is that residual error is
decomposed rather than reported as a single aggregate figure. Three
contributions are independent in origin and demand different remedies;
conflating them would make the aggregate uninformative about what to improve.

\begin{enumerate}[leftmargin=2em]
  \item \emph{Landmark detection noise.} Error in the image-plane hand
        landmarks, propagated through depth lifting. This contribution is
        reduced by confidence weighting \eqref{eq:weights} and by multi-view
        fusion, and it is characterised by the wrist-position repeatability
        measurement of Table~\ref{tab:calib}. It is not reduced by better
        calibration.
  \item \emph{Calibration error.} Error in the camera-to-workspace and
        robot-to-workspace transforms \eqref{eq:extrinsic}. This contribution
        is systematic within a session, appears as a constant offset rather
        than as noise, and is characterised by the reprojection and
        cross-camera registration measurements. It is not reduced by
        multi-view fusion, since all views inherit it.
  \item \emph{Representation dependency.} Error introduced by the choice of
        reference frame, specifically the sensitivity of world-frame
        trajectories to object displacement. This contribution is zero when
        the object is where it was during demonstration and grows with
        displacement; it is characterised by the comparison of
        Table~\ref{tab:repr} rather than by any calibration measurement.
\end{enumerate}

\begin{table}[htbp]
\caption{Decomposition of Residual Wrist-Pose Error by Source}
\label{tab:errbudget}
\centering
\small
\begin{tabular}{llcc}
\toprule
Source & Character & Magnitude (mm) & Reduced by \\
\midrule
Landmark detection & random  & \ph{X.X} & fusion, confidence weighting \\
Calibration        & systematic & \ph{X.X} & board size, view count \\
Representation     & conditional & \ph{X.X} & object-relative frame \\
\midrule
Combined           & --- & \ph{XX.X} & --- \\
\bottomrule
\end{tabular}
\end{table}

\TODO{Populate and discuss. If the combined figure differs materially from the
quadrature sum of the independent contributions, the sources are not
independent and the decomposition must be revised rather than presented as if
it held.}

\section{Limitations}
\label{sec:limitations}

The assumptions of Section~\ref{sec:assumptions} translate into the following
limitations, which are stated plainly rather than deferred to future work.

\begin{enumerate}[leftmargin=2em]
  \item \emph{Gripper abstraction.} A binary gripper state cannot represent
        in-hand manipulation, regrasping or force-modulated grasps. Tasks
        requiring these are outside the demonstrated capability, and extending
        to dexterous hands would require reinstating the finger-level
        representation that Section~\ref{sec:lit-hand} rejects.
  \item \emph{Rigid-object assumption.} The object-relative transform
        \eqref{eq:objrel} is undefined for deformable objects. Cloth, cable and
        granular manipulation are excluded.
  \item \emph{Dependence on object pose estimation.} The object-relative
        representation inherits the error of the object pose estimator. Where
        that estimator fails---textureless, symmetric or occluded objects---the
        representation degrades to worse than the world-frame alternative it is
        meant to improve upon.
  \item \emph{Kinematic transfer only.} No force or contact model is estimated.
        Success on tasks with a force-regulation component is not claimed.
  \item \emph{Single demonstrator, constrained task class.} The evaluation uses
        \ph{one} demonstrator and \ph{three} task types. Generalisation across
        demonstrators, hand morphologies and task families is not established.
  \item \emph{Sensing floor.} Cross-camera agreement at the centimetre level
        \cite{romeo} bounds achievable accuracy. Tasks with tolerances below
        that bound cannot be addressed by this sensing configuration
        regardless of the retargeting formulation.
  \item \emph{Cross-family synchronisation.} Alignment between the Kinect pair
        and the RealSense is software-based and therefore millisecond-scale;
        fast motion may exhibit residual misalignment between those streams.
\end{enumerate}

\section{Threats to Validity}
\label{sec:validity}

\paragraph{Internal validity.} Metrics were fixed before the final runs
(Section~\ref{sec:metrics}) to prevent selection of criteria after observing
outcomes. Baselines B1--B5 differ from the proposed method in exactly one
component each, so attribution of any observed difference is unambiguous. The
principal residual threat is that the sequential baseline's filter parameters
are themselves tunable; an unfavourable choice would make the comparison
misleading. \TODO{State how B1 filter parameters were selected---ideally by
tuning them to optimise the baseline's own fidelity, not by adopting a default.}

\paragraph{External validity.} The constrained task class was chosen for
measurability, which necessarily limits generalisation. Results should not be
read as claims about manipulation in general.

\paragraph{Construct validity.} Spectral arc length measures smoothness of the
velocity profile, which is a proxy for the property actually of interest---that
a policy can imitate the trajectory. The proxy is standard but imperfect, and
the downstream policy results of Section~\ref{sec:res-policy} are the more
direct evidence.

\section{Implications}
\label{sec:implications}

Three implications extend beyond the specific system.

First, the transferability of the online retargeting formulation to the offline
regime suggests that the separation between the teleoperation and
video-imitation literatures is a historical accident rather than a technical
necessity. The offline setting is strictly less constrained---no causality
requirement, no latency budget---and therefore able to use the same formulation
more thoroughly than the setting for which it was designed.

Second, moving curation upstream of export, so that infeasible trajectories are
removed before they enter the dataset rather than filtered after they have
degraded a policy, is applicable to any pipeline that synthesises rather than
records actions. The cost is one physical execution per episode; the benefit is
that the exported dataset contains proprioception the robot actually attained.

Third, the error decomposition of Section~\ref{sec:errbudget} is a
transferable methodological point. Reporting a single aggregate accuracy figure
for a multi-stage perception pipeline obscures which stage limits performance,
and consequently which engineering effort would pay.

\chapter{Conclusion}
\label{ch:concl}

\section{Summary}
\label{sec:summary}

This thesis has addressed the question of how manipulation demonstrations
recorded passively, with a human using their own hands and no teleoperation
interface, can be converted into datasets suitable for training visuomotor
imitation policies. The work has argued that the two established approaches to
this problem are each incomplete: online teleoperation systems have converged
on a sound retargeting formulation---a single weighted cost functional
combining data fidelity, temporal smoothness and kinematic constraints---but
produce no dataset, while offline video-to-robot pipelines produce datasets but
retarget either not at all, through sparse object keyframes, or through a
sequential filter--solve--filter pattern whose smoothing is applied in the wrong
space to guarantee the property it is intended to provide.

\viki\ has been presented as a system occupying the intersection. It captures
demonstrations with a metrically calibrated multi-camera RGB-D rig anchored to
a ChArUco workspace frame; represents the hand minimally, as a wrist pose in
$\SE$ with a binary gripper state; stores each demonstration both relative to
the manipulated object and relative to the workspace anchor; retargets by
minimising one weighted functional in which the data term is scaled by
per-frame perception confidence and robustified by a Huber loss, and in which
velocity and acceleration penalties act directly in joint space; validates the
result by executing it on the physical manipulator before export; and writes a
LeRobot-compatible dataset containing the proprioception the robot actually
attained rather than the states the optimiser predicted.

\section{Answers to the Research Questions}
\label{sec:answers}

\paragraph{RQ1.} \ph{[State the measured effect of the single-functional
formulation on fidelity and smoothness relative to the sequential baseline, and
whether the trade-off frontier moved rather than merely shifted along.]}

\paragraph{RQ2.} \ph{[State the achieved wrist-pose accuracy and how the
residual decomposes across landmark noise, calibration error and
representation dependency, per Table~\ref{tab:errbudget}.]}

\paragraph{RQ3.} \ph{[State the proportion of episodes rejected by replay
screening, the dominant rejection causes, and the resulting change in the
executable fraction of the exported dataset.]}

\section{Contributions}
\label{sec:contrib-final}

The contributions of this work are the transfer of the single-cost-functional
retargeting formulation from the online teleoperation regime, for which it was
developed, into the offline object-centric dataset-generation regime, where the
review of Chapter~\ref{ch:lit} found it unused; the coupling of that
formulation to a metrically calibrated multi-view RGB-D front end with
confidence-weighted fusion and explicit cross-device timestamp handling; the
use of physical replay as a pre-export dataset-curation stage rather than a
post-training evaluation; and an error decomposition that separates the three
independent contributions to residual pose error rather than reporting an
aggregate.

It is worth restating what is \emph{not} claimed. The cost functional is not
novel as an operator; its components appear in \cite{anyteleop, dexflow,
opentv, dexteleop0}, and a claim of algorithmic novelty would be refuted by
that prior art. Object-centric representation, ChArUco anchoring,
confidence-weighted depth fusion and host-clock synchronisation are all
established practice. The contribution is the conjunction and the regime
transfer, which the positioning analysis of Section~\ref{sec:positioning} found
unoccupied.

\section{Limitations}
\label{sec:concl-limits}

The system is restricted to parallel-jaw targets, rigid objects, single-arm
quasi-static tabletop tasks and a fixed workspace; it performs kinematic
transfer only, estimating no forces; its accuracy is bounded below by
centimetre-level cross-camera agreement; and its evaluation uses a single
demonstrator on a constrained task class. These are stated in full in
Section~\ref{sec:limitations}.

\section{Future Work}
\label{sec:future}

Several directions follow from the limitations.

\begin{enumerate}[leftmargin=2em]
  \item \emph{Dexterous targets.} Extending beyond a binary gripper requires
        reinstating a finger-level representation. The ablation evidence of
        \cite{keyobj} suggests that the wrist objective should remain weakly
        weighted relative to fingertip objectives, which would change the
        balance of terms in \eqref{eq:cost} rather than its structure.
  \item \emph{Deformable objects.} Replacing the rigid object pose of
        \eqref{eq:objrel} with a flow-based object representation
        \cite{novaflow} would remove the rigidity assumption at the cost of a
        less compact transfer rule.
  \item \emph{Closing the replay loop.} Currently, replay failures trigger
        bounded re-solution with hand-adjusted weights. Learning the weight
        adjustment from the failure mode would make the loop automatic.
  \item \emph{Cross-embodiment transfer.} The object-relative representation is
        embodiment-agnostic by construction; the same recorded demonstrations
        could be retargeted to the KUKA iiwa14 available in the laboratory, and
        the degree to which policy performance transfers would test the
        representational claim directly.
  \item \emph{Scaling demonstrator diversity.} Recording from multiple
        demonstrators with differing hand morphologies would test whether the
        fixed hand-to-end-effector offset $\bT^{H}_{E}$ of
        \eqref{eq:transfer} requires per-demonstrator calibration.
  \item \emph{Downstream curation.} The confidence weights and replay residuals
        retained in the exported dataset support curation experiments in the
        manner of \cite{demoscore} without re-recording; whether the replay
        residual is a better curation signal than a learned rollout classifier
        is an open question.
\end{enumerate}

\section{Concluding Remarks}
\label{sec:remarks}

The practical claim of this work is modest and, for that reason, defensible:
a formulation already accepted in one part of the robot-learning literature
works in another part where it had not been applied, and the infrastructure
required to apply it---metric multi-view sensing, honest cross-device
synchronisation, physical validation before export---is buildable at the scale
of a single laboratory. If the results hold, the cost of producing
imitation-learning data falls without the cost being transferred to policy
robustness, which is the trade that passive capture has historically demanded.

\chapter*{References}
\addcontentsline{toc}{chapter}{References}

\begin{thebibliography}{99}
\setlength{\itemsep}{2pt}

\bibitem{anyteleop}
Y. Qin, W. Yang, B. Huang, K. Van Wyk, H. Su, X. Wang, Y.-W. Chao, and D. Fox,
``AnyTeleop: A general vision-based dexterous robot arm-hand teleoperation
system,'' in \emph{Proc. Robotics: Science and Systems (RSS)}, 2023.
[Online]. Available: \url{https://arxiv.org/abs/2307.04577}

\bibitem{opentv}
X. Cheng, J. Li, S. Yang, G. Yang, and X. Wang, ``Open-TeleVision:
Teleoperation with immersive active visual feedback,'' in \emph{Proc. Conf. on
Robot Learning (CoRL)}, 2024. [Online]. Available:
\url{https://arxiv.org/abs/2407.01512}

\bibitem{bunny}
R. Ding, Y. Qin, J. Zhu, C. Jia, S. Yang, R. Yang, X. Qi, and X. Wang,
``Bunny-VisionPro: Real-time bimanual dexterous teleoperation for imitation
learning,'' in \emph{Proc. IEEE/RSJ Int. Conf. on Intelligent Robots and
Systems (IROS)}, 2025. [Online]. Available:
\url{https://arxiv.org/abs/2407.03162}

\bibitem{ace}
S. Yang, M. Liu, Y. Qin, R. Ding, J. Li, X. Cheng, R. Yang, S. Yi, and X. Wang,
``ACE: A cross-platform visual-exoskeletons system for low-cost dexterous
teleoperation,'' in \emph{Proc. Conf. on Robot Learning (CoRL)}, 2024.
[Online]. Available: \url{https://arxiv.org/abs/2408.11805}

\bibitem{dexflow}
\ph{[Authors]}, ``DexFlow: Dexterous retargeting with $C^2$ continuity via
Hessian regularisation,'' \emph{arXiv:2505.01083}, 2025. \ph{[Verify author
list and venue.]}

\bibitem{dexteleop0}
\ph{[Authors]}, ``\ph{[Title]},'' 2025. \ph{[Robust Huber-loss retargeting
formulation cited in Section~\ref{sec:lit-retarget}; confirm identifier before
submission.]}

\bibitem{keyobj}
\ph{[C.]} Xin, \ph{[Y.]} Yu, \ph{[X.]} Jiang, \ph{[Y.]} Zhang, and \ph{[H.]}
Li, ``Analyzing key objectives in human-to-robot retargeting for dexterous
manipulation,'' \emph{arXiv:2506.09384}, 2025.

\bibitem{pink}
S. Caron \emph{et al.}, ``Pink: Python inverse kinematics based on Pinocchio,''
software library. [Online]. Available:
\url{https://github.com/stephane-caron/pink}

\bibitem{orion}
Y. Zhu, \ph{[et al.]}, ``Vision-based manipulation from single human video with
open-world object graphs,'' \emph{arXiv:2405.20321}, 2024; \emph{Autonomous
Robots}, 2026. \ph{[Confirm final volume and pages.]}

\bibitem{okami}
J. Li, Y. Zhu, Y. Xie, Z. Jiang, M. Seo, G. Pavlakos, and Y. Zhu, ``OKAMI:
Teaching humanoid robots manipulation skills through single video imitation,''
in \emph{Proc. Conf. on Robot Learning (CoRL)}, \emph{PMLR}, vol.~270,
pp.~299--317, 2024.

\bibitem{dexmv}
Y. Qin, Y.-H. Wu, S. Liu, H. Jiang, R. Yang, Y. Fu, and X. Wang, ``DexMV:
Imitation learning for dexterous manipulation from human videos,'' in
\emph{Proc. European Conf. on Computer Vision (ECCV)}, 2022.

\bibitem{spot}
\ph{[Authors]}, ``SPOT: SE(3) pose trajectory diffusion for object-centric
manipulation,'' in \emph{Proc. IEEE Int. Conf. on Robotics and Automation
(ICRA)}, 2025. [Online]. Available: \url{https://arxiv.org/abs/2411.00965}

\bibitem{novaflow}
\ph{[Authors]}, ``NovaFlow: Zero-shot manipulation via actionable flow from
generated videos,'' \emph{arXiv:2510.08568}, 2025.

\bibitem{demoscore}
A. Chen, A. Lessing, Y. Liu, and C. Finn, ``Curating demonstrations using
online experience,'' \emph{arXiv:2503.03707}, 2025.

\bibitem{openvlaoft}
M. J. Kim, C. Finn, and P. Liang, ``Fine-tuning vision-language-action models:
Optimizing speed and success,'' \emph{arXiv:2502.19645}, 2025.

\bibitem{egomimic}
S. Kareer, D. Patel, R. Punamiya, P. Mathur, S. Cheng, C. Wang, J. Hoffman, and
D. Xu, ``EgoMimic: Scaling imitation learning via egocentric video,'' in
\emph{Proc. Conf. on Robot Learning (CoRL)}, 2024. [Online]. Available:
\url{https://arxiv.org/abs/2410.24221}

\bibitem{dexwild}
T. Tao, M. K. Srirama, J. J. Liu, K. Shaw, and D. Pathak, ``DexWild: Dexterous
human interactions for in-the-wild robot policies,'' in \emph{Proc. Robotics:
Science and Systems (RSS)}, 2025. [Online]. Available:
\url{https://arxiv.org/abs/2505.07813}

\bibitem{toolif}
\ph{[Authors]}, ``Tool-as-Interface: Learning robot policies from observing
human tool use,'' \emph{arXiv:2504.04612}, 2025.

\bibitem{howtransfer}
\ph{[Authors]}, ``Hand-centric human-to-robot trajectory transfer from video
demonstrations via open-world contact localization,'' \emph{arXiv:2606.10743},
2026. \ph{[Preprint; peer-review status unconfirmed.]}

\bibitem{viroil}
\ph{[Authors]}, ``Video-to-robot imitation learning for bimanual dexterous
manipulation from human video,'' in \emph{Proc. ICEILM}, 2025,
doi:10.1145/3788731.3788736. \ph{[Verify author list from the published PDF.]}

\bibitem{robopaint}
\ph{[Authors]}, ``RoboPaint: From human demonstration to any robot and any
view,'' \emph{arXiv:2602.05325}, 2026. \ph{[Preprint; status unconfirmed.]}

\bibitem{shadowing}
\ph{[Authors]}, ``Vision-based hand shadowing,'' \emph{IEEE Access}, 2026;
\emph{arXiv:2603.11383}. \ph{[Preprint; status unconfirmed.]}

\bibitem{dexcanvas}
\ph{[Authors]}, ``DexCanvas: Bridging human demonstrations and robot learning
for dexterous manipulation,'' \emph{arXiv:2510.15786}, 2025.

\bibitem{romeo}
L. Romeo, R. Marani, A. G. Perri, and T. D'Orazio, ``Microsoft Azure Kinect
calibration for three-dimensional dense point clouds and reliable skeletons,''
\emph{Sensors}, vol.~22, no.~13, art.~4986, 2022.

\bibitem{darwish}
W. Darwish, Q. Bols\'{e}e, and A. Munteanu, ``Robust calibration of a
multi-view Azure Kinect scanner based on spatial consistency,'' in
\emph{Proc. IEEE Int. Conf. on 3D Immersion (IC3D)}, 2020.

\bibitem{kimcal}
J. Kim \emph{et al.}, ``Robust extrinsic calibration of multiple RGB-D cameras
with body tracking and feature matching,'' \emph{Sensors}, vol.~21, no.~3,
art.~1013, 2021.

\bibitem{bulee}
S. Bu and S. Lee, ``Easy to calibrate: Marker-less calibration of multiview
Azure Kinect,'' \emph{Computer Modeling in Engineering \& Sciences},
vol.~136, no.~3, pp.~3083--3096, 2023.

\bibitem{tolgyessy}
M. T\"{o}lgyessy, M. Dekan, \v{L}. Chovanec, and P. Hubinsk\'{y}, ``Evaluation
of the Azure Kinect and its comparison to Kinect v1 and Kinect v2,''
\emph{Sensors}, vol.~21, no.~2, art.~413, 2021.

\bibitem{kurillo}
G. Kurillo, E. Hemingway, M.-L. Cheng, and L. Cheng, ``Evaluating the accuracy
of the Azure Kinect and Kinect v2,'' \emph{Sensors}, vol.~22, no.~7, art.~2469,
2022.

\bibitem{mskinect}
Microsoft, ``Azure Kinect DK hardware specification'' and ``Synchronize
multiple Azure Kinect DK devices,'' Microsoft Learn documentation.
\ph{[Add access date.]}

\bibitem{intel}
Intel Corporation, ``Intel RealSense D400 series product family datasheet'' and
multi-camera configuration guidance. \ph{[Add revision and access date.]}

\bibitem{ridgerun}
RidgeRun, ``Multi-camera interface configurations and synchronization,''
technical note. \ph{[Non-peer-reviewed vendor documentation; retain only for
engineering practice, not for scientific claims.]}

\bibitem{alihandeye}
I. Ali \emph{et al.}, ``Accuracy evaluation of hand-eye calibration techniques
for vision-guided robots,'' \emph{PLOS ONE}, 2022. \ph{[Add volume and article
number.]}

\bibitem{fusionpatent}
``Image-based geometric fusion of multiple depth images using ray casting,''
U.S. Patent 10\,818\,071, 2020.

\bibitem{zeng}
J. Zeng, Y. Tong, Y. Huang, Q. Yan, W. Sun, J. Chen, and Y. Wang, ``Deep
surface normal estimation with hierarchical RGB-D fusion,'' in \emph{Proc. IEEE
Conf. on Computer Vision and Pattern Recognition (CVPR)}, 2019. [Online].
Available: \url{https://arxiv.org/abs/1904.03405}

\bibitem{hrbf}
Y. Wang \emph{et al.}, ``HRBF-Fusion: Accurate 3D reconstruction from RGB-D
data using on-the-fly implicits,'' \emph{ACM Transactions on Graphics}.
\ph{[Add volume, issue and year.]}

\end{thebibliography}

\vspace{4mm}
\noindent\ph{\small Note: entries marked in this colour require completion or
verification before submission. Preprint entries dated 2026 could not be
confirmed as peer-reviewed at the time of writing and are cited as preprints.}

\appendix
\chapter{Notation}
\label{app:notation}

\begin{table}[htbp]
\caption{Symbols Used Throughout the Thesis}
\label{tab:notation}
\centering
\small
\begin{tabular}{ll}
\toprule
Symbol & Meaning \\
\midrule
$W$ & workspace frame, defined by the ChArUco anchor \\
$C_k$ & frame of camera $k$ \\
$R$ & manipulator base frame \\
$E$ & manipulator end-effector frame \\
$H$ & human wrist frame \\
$O$ & manipulated object frame \\
$\bT^{A}_{B}$ & rigid transform from frame $B$ to frame $A$, element of $\SE$ \\
$\bq_t \in \R^{n}$ & manipulator configuration at timestep $t$ \\
$f(\bq_t)$ & forward kinematics of the end-effector \\
$\be_t$ & task-space residual, \eqref{eq:residual} \\
$\omega_t$ & per-frame perception confidence weight, \eqref{eq:conf} \\
$w^{k}_{i,t}$ & per-landmark fusion weight, \eqref{eq:weights} \\
$v^{k}_{i,t}$ & detector visibility score \\
$c^{k}_{i,t}$ & sensor depth confidence \\
$\rho_\delta$ & Huber loss with transition parameter $\delta$ \\
$\lambda_v,\lambda_a,\lambda_r$ & velocity, acceleration and posture weights \\
$g_t$ & binary gripper state, \eqref{eq:gripper} \\
$\tau$ & host-side monotonic timestamp, \eqref{eq:hostclock} \\
$T$ & number of timesteps in an episode \\
\bottomrule
\end{tabular}
\end{table}

\chapter{Exported Dataset Schema}
\label{app:schema}

Each episode is stored as an HDF5 group with the fields listed in
Table~\ref{tab:schema}. The layout is compatible with LeRobot conventions;
fields beyond the LeRobot minimum are retained so that curation can be repeated
downstream without re-recording.

\begin{table}[htbp]
\caption{Per-Episode Dataset Fields}
\label{tab:schema}
\centering
\small
\begin{tabular}{lll}
\toprule
Field & Shape & Description \\
\midrule
\texttt{observation.images.*} & $T\times H\times W\times 3$ & per-camera colour
  streams \\
\texttt{observation.state} & $T\times n$ & recorded joint positions from replay \\
\texttt{action} & $T\times (n{+}1)$ & commanded joint positions and gripper \\
\texttt{wrist\_pose\_world} & $T\times 4\times 4$ & $\bT^{W}_{H,t}$ \\
\texttt{wrist\_pose\_object} & $T\times 4\times 4$ & $\bT^{O}_{H,t}$,
  \eqref{eq:objrel} \\
\texttt{object\_pose\_world} & $T\times 4\times 4$ & $\bT^{W}_{O,t}$ \\
\texttt{gripper\_state} & $T$ & binary, \eqref{eq:gripper} \\
\texttt{confidence} & $T$ & $\omega_t$, \eqref{eq:conf} \\
\texttt{replay\_residual} & $T$ & controller tracking residual during replay \\
\texttt{timestamp\_host} & $T$ & host monotonic timestamp, \eqref{eq:hostclock} \\
\bottomrule
\end{tabular}
\end{table}

\chapter{Solver Configuration}
\label{app:config}

\begin{table}[htbp]
\caption{Weights and Solver Parameters Used in the Reported Experiments}
\label{tab:config}
\centering
\small
\begin{tabular}{lll}
\toprule
Parameter & Value & Selection method \\
\midrule
$\lambda_v$ (velocity)      & \ph{X.XX} & \ph{[grid search / held-out set]} \\
$\lambda_a$ (acceleration)  & \ph{X.XX} & \ph{[\dots]} \\
$\lambda_r$ (posture)       & \ph{X.XX} & \ph{[\dots]} \\
$\alpha$ (confidence sharpness) & \ph{X.XX} & \ph{[\dots]} \\
$\delta$ (Huber transition) & \ph{X.XX} & \ph{[\dots]} \\
$\tau_g$ (gripper threshold, mm) & \ph{XX} & \ph{[\dots]} \\
LM damping                  & \ph{X.XX} & \ph{[\dots]} \\
Recording rate (Hz)         & \ph{XX}   & --- \\
Max replay re-solutions     & \ph{X}    & --- \\
\bottomrule
\end{tabular}
\end{table}

\TODO{Weight selection must be described honestly: if weights were tuned on the
same episodes used for evaluation, say so and report the consequence, or
re-tune on a held-out split.}


\end{document}
